\emph{Översikt}
Vad kan du nu, och hur hänger det ihop?  Det här kapitlet repeterar hela
kursen, ett ämne i taget och i samma ordning som föreläsningarna: från
algoritmer, variabler och funktioner, via inmatning, villkor,
upprepningar och behållare, till klasser, filer och grafiska
gränssnitt.  Varje avsnitt börjar med en fråga, och svaret är ett steg i
samma program, en bank: först ett saldo i en variabel, sist klasser med
ett fönster framför.  Efter banken kommer konstruktionens form i en
syntaxruta, ett kort exempel i ett annat sammanhang och var
konstruktionen gicks igenom första gången.  Kapitlet slutar med det som
lådorna i bilderna har förenklat: egentligen är alla namn pilar.

\emph{Lärandemål}
Efter denna modul ska studenten kunna:
% Lo04 (styrstrukturer), Lo09 (sammansatta datatyper), Lo10 (filer)
\begin{restatable}{lo}{RecapLOConstructs}\label{RecapLOConstructs}%
  Känna igen och skriva kursens konstruktioner utifrån deras form:
  tilldelning, funktioner, villkor, upprepningar, behållare, klasser,
  filer och fönster.
\end{restatable}
% Lo11 (granska program)
\begin{restatable}{lo}{RecapLOPredict}\label{RecapLOPredict}%
  Förutsäga vad ett kort program skriver ut, och förklara varför med
  hjälp av en bild av namnen och värdena.
\end{restatable}
% Lo02 (dela upp problem), Lo03 (dela upp program)
\begin{restatable}{lo}{RecapLOConnect}\label{RecapLOConnect}%
  Följa ett och samma problem genom kursens konstruktioner, och förklara
  vad varje ny konstruktion löser som den förra inte gjorde.
\end{restatable}
% Lo09 (sammansatta datatyper)
\begin{restatable}{lo}{RecapLOChoose}\label{RecapLOChoose}%
  Välja konstruktion efter problemet: funktion, villkor eller upprepning,
  och lista, tupel, uppslagslista eller klass.
\end{restatable}
% Lo09 (sammansatta datatyper)
\begin{restatable}{lo}{RecapLONames}\label{RecapLONames}%
  Avgöra om en ändring syns hos den som anropade en funktion, genom att
  skilja på att flytta ett namn och att ändra ett objekt.
\end{restatable}

\emph{Förkunskaper}
Repetitionen bygger på alla kursens föreläsningar, från
\emph{Algoritmiskt tänkande} till \emph{Grafiska användargränssnitt}.  Den
som ännu inte har tagit del av de sista --- från \emph{Behållare:
Mängder, stackar och köer} och framåt --- kan läsa de avsnitten som en
förhandstitt på vad som kommer.
